\subsection{应用：代数的自同构群的不变量}

给定环 K、K-代数 A 和群 G，若满足：(1) 集 A 以 G 为算子群（《代数》，第 I 章，§ 7，第 2 节）；(2) 对所有 \(\sigma \in G\) ，映射 \(x \mapsto \sigma \cdot x\) 是 K-代数 A 的自同态（且由于它是双射（同前），因而是自同构），则我们说 G 作用于 A。我们将以 \(A^{G}\) 表示 A 中在 G 下不变的元素组成的集；显然它是 A 的子 K-代数。

若 \(\mathcal{G}\) 在 A 中的每条轨道（《代数》，第 I 章，第 IV 分册勘误）都是有限的，则我们说 \(\mathcal{G}\) 是 A 上局部有限的算子群。

命题 22. 设 A 是（交换）K-代数，G 是 A 上局部有限的算子群。则 A 在子代数 \(\mathbf{A}^{\mathcal{G}}\) 上是整的。

对所有 \(x \in \mathbf{A}\) ，设 \(x_{i} (1 \leqslant i \leqslant n)\) 是 \(x\) 在 \(\mathcal{G}\) 下轨道中的不同元素；对所有 \(\sigma \in \mathcal{G}\) ，存在置换 \(\pi_{\sigma}\) （集 \(\{1, 2, \ldots, n\}\) 的置换），使得 \(\sigma. x_{i} = x_{\pi_{\sigma}(i)}\) （\(1 \leqslant i \leqslant n\)）；因此 \(x_{i}\) 的初等对称函数是 \(\mathbf{A}\) 中在 \(\mathcal{G}\) 下不变的元素，换言

之，是 \(\mathbf{A}^{\mathcal{G}}\) 的元素。由于 \(x\) 是首一多项式 \(\prod_{i=1}^{n} (\mathbf{X} - x_i)\) 的根，且该多项式的系数属于 \(\mathbf{A}^{\mathcal{G}}\) ，故 \(x\) 在 \(\mathbf{A}^{\mathcal{G}}\) 上是整的。

定理 2. 设 A 是有限生成的 K-代数，G 是 A 上局部有限的算子群。则 A 是有限生成的 \(\mathbf{A}^{\mathcal{G}}\)-模；若进一步 K 是 Noether 环，则 \(A^{G}\) 是有限生成的 K-代数。

设 \((a_{j})_{1\leqslant j\leqslant m}\) 是 K-代数 A 的生成系；由于更不必说 \(\mathbf{A} = \mathbf{A}^{\mathcal{G}}[a_1,\ldots ,a_m]\) ，且由命题 22 诸 \(a_{j}\) 在 \(\mathbf{A}^{\mathcal{G}}\) 上是整的，故第一个断言由第 1 节命题 4 得出。第二个断言是下述引理的推论：

引理 5. 设 K 是 Noether 环，B 是有限生成的 K-代数，C 是 B 的子 K-代数且 B 在 C 上是整的。则 C 是有限生成的 K-代数。

设 \((x_{i})_{1\leqslant i\leqslant n}\) 是 K-代数 B 的有限生成系。对所有 \(i\) ，按假设存在首一多项式 \(\mathbf{P}_i\in \mathbf{C}[\mathbf{X}]\) 使得 \(\mathrm{P_i}(x_i) = 0\) 。设 \(\mathbf{C}'\) 是由诸 \(\mathbf{P}_i\) （ \(1\leqslant i\leqslant n\) ）的系数生成的 C 的子 K-代数；显然诸 \(x_{i}\) 在 \(\mathbf{C}'\) 上是整的，且 \(\mathbf{B} = \mathbf{C}'[x_1,\dots,x_n]\) ；因此 B 是有限生成的 \(\mathbf{C}'\) -模（第 1 节，命题 4）。另一方面，\(\mathbf{C}'\) 是 Noether 环（第 III 章，§2，第 10 节，定理 2 的推论 3）；因此 C 是有限生成的 \(\mathbf{C}'\) -模，这就证明了 C 是有限生成的 K-代数。

注. 由所有 \(\sigma \in \mathcal{G}\) （满足 \(\sigma a_{j} = a_{j}\) ，\(1 \leqslant j \leqslant m\) ）组成的集显然使 A 的每个元素不变。因此，使 A 的每个元素都不变的正规子群 \(\mathcal{H}\) （\(\mathcal{G}\) 的）在 \(\mathcal{G}\) 中的指数有限，从而可以把 A 看作以有限群 \(\mathcal{G}/\mathcal{H}\) 为算子群；显然 \(\mathrm{A}^{\mathcal{G}/\mathcal{H}} = \mathrm{A}^{\mathcal{G}}\) 。

设 S 是环 A 的乘性子集，G 是作用于 A 且使 S 稳定的群；则对所有 \(\sigma\in G\) ，存在唯一的自同态 \(z\mapsto\sigma.z\) （环 \(S^{-1}A\) 的），使得 \(\sigma.(a/1)=(\sigma.a)/1\) （对所有 \(a\in A\)）；它由公式 \(\sigma.(a/s)=(\sigma.a)/( \sigma.s)\) 给出，其中 \(a\in A\) ，\(s\in S\) （第 II 章，§2，第 1 节，命题 2）；若 \(\tau\) 是 G 的另一个元素，则显然 \(\sigma.(\tau.z)=(\sigma\tau).z\) 对所有 \(z\in S^{-1}A\) 成立，因而群 G 作用于环 \(S^{-1}A\) 。

命题 23. 设 A 是 K-代数，\(\mathcal{G}\) 是 A 上局部有限的算子群，S 是 A 的在 \(\mathcal{G}\) 下稳定的乘性子集，\(\mathrm{S}^{\mathcal{G}}\) 是集 \(\mathrm{S} \cap \mathrm{A}^{\mathcal{G}}\) 。则从 \((\mathrm{S}^{\mathcal{G}})^{-1}\mathrm{A}\) 到 \(\mathrm{S}^{-1}\mathrm{A}\) 的典范映射（第 II 章，§ 2，第 1 节，命题 2 的推论 2）是同构，且它把 \((\mathrm{S}^{\mathcal{G}})^{-1}\mathrm{A}^{\mathcal{G}}\) 映到 \((\mathrm{S}^{-1}\mathrm{A})^{\mathcal{G}}\) 。

对所有 \(s \in S\) ，设 \(s, s_1, \ldots, s_q\) 是 \(s\) 在 \(\mathcal{G}\) 下轨道中的不同元素；由于 \(ss_1 \ldots s_q \in S^{\mathcal{G}}\) ，第一个断言由第 II 章，§ 2，第 3 节，命题 8 得出。把 \((S^{\mathcal{G}})^{-1}A\) 与 \(S^{-1}A\) 典范地等同起来，则显然 \((S^{\mathcal{G}})^{-1}A^{\mathcal{G}}\) 的每个元素都在 \(\mathcal{G}\) 下不变。反之，设 \(a/t\) 是 \((S^{\mathcal{G}})^{-1}A\) 中在 \(\mathcal{G}\) 下不变的元素（ \(a \in A\) ，\(t \in S^{\mathcal{G}}\) ）；若 \(a_j\) （ \(1 \leqslant j \leqslant m\) ）是 \(a\) 在 \(\mathcal{G}\) 下轨道中的不同元素，则 \(a_j/t = a/t\) （\(1 \leqslant j \leqslant m\)），因此存在 \(s \in S^{\mathcal{G}}\) 使得 \(s(a_j - a) = 0\) （\(1 \leqslant j \leqslant m\)）；换言之，\(sa\) 在 \(\mathcal{G}\) 下不变，且由于 \(a/t = (sa)/(st)\) ，当然有 \(a/t \in (S^{\mathcal{G}})^{-1}A^{\mathcal{G}}\) 。

推论. 设 A 是整环，K 是其分式域，G 是 A 上局部有限的算子群。则 G 作用于 K，且 \(K^{G}\) 是 \(A^{G}\) 的分式域。

A - \{0\} 在 G 下稳定。

